\clearpage\begingroup
\NoSourceHead\SourceLeaf{125}
\begin{center}
[II.]\par\medskip
{\large [ANALYSE VON FERNROHREN.]}\par\medskip
\rule{.55\linewidth}{.4pt}\par
{\small [Aus Handbuch 20, Bf, S. 306, 308, 309.]}\par
\rule{.55\linewidth}{.4pt}\par\bigskip
[A.]\par\medskip
\textbf{Dimensionen des Baumannschen Kometensuchers.}
\end{center}
\begin{center}\small
\setlength{\tabcolsep}{5pt}\renewcommand{\arraystretch}{1.45}
\begin{tabular}{l|r|c|c|c}
\hline\hline
{}[Linsen]&Brennweite&Durchmesser&Dicke&Distanzen\\
\hline\hline
Objectiv&$862$&$52$&$15$&$772{,}7[-x=h]$\\
{}[Okular] 1&$5576=181{,}8$&$60{,}8$&$6{,}5$&\raisebox{-.5\normalbaselineskip}[0pt][0pt]{$94{,}2\ [=h']$}\\
\hfill 2&$2984=97{,}3$&$45{,}4$&$5{,}4$&\raisebox{-.5\normalbaselineskip}[0pt][0pt]{$14{,}2\ [=h'']$}\\
\hfill 3&$2065=67{,}3$&$38$&$5{,}0$&
\end{tabular}
\end{center}
\noindent\rule{\linewidth}{.4pt}
\[
\begin{aligned}
[f\ ={}]&\ 1:862\quad[=0{,}00116009]\\
{}[f'\ ={}]&\ 1:181{,}8\ [=0{,}00550055]\\
{}[f''\ ={}]&\ 1:\phantom{1}97{,}3\ [=0{,}01027749]\\
{}[f'''\ ={}]&\ 1:\phantom{1}67{,}3\ [=0{,}01485884]
\end{aligned}
\]
\noindent\rule{\linewidth}{.4pt}
\[
\begin{aligned}
\alpha&=+1\\
{}[(f,h)\ ={}]\ \alpha'&=-0{,}103596[-0{,}00116009x]\\
{}[(f,h')\ ={}]\ \alpha''&=-0{,}059363[+0{,}00055899x]\\
{}[(f,h'')\ ={}]\ \alpha'''&=+0{,}075266[-0{,}00038681x]
\end{aligned}
\]

\clearpage\SourceLeaf{126}\HistoricalHead{121}{R}{ANALYSE VON FERNROHREN.}\HistoricalFoot{XI$_1$.}{16}
\[
\begin{aligned}
[f={}]\ \beta&=+0{,}00116009\\
{}[(f,f')\ ={}]\ \beta'&=-0{,}00172993[-0{,}00000638x]\\
{}[(f,f'')\ ={}]\ \beta''&=+0{,}00111983[+0{,}0000121263x]\\
{}[(f,f''')\ ={}]\ \beta'''&=-0{,}00000146[-0{,}0000178676x]
\end{aligned}
\]
\noindent\rule{\linewidth}{.4pt}
\[
\begin{aligned}
\gamma&=0\\
\gamma'&=+772{,}7\ [-1{,}0000x]\\
\gamma''&=-466{,}52[+0{,}48185x]\\
\gamma'''&=+352{,}28[-0{,}333219x]\text{[*)]}
\end{aligned}
\]
\begingroup\renewcommand{\thefootnote}{}\footnotetext{[*) In der Handschrift befindet sich hier ein Rechenfehler (342,28 statt 352,28), der auch $\delta'''$ und den Wert von $x$ beeinflußt hat; im Folgenden sind die richtigen Werte von $\delta'''$ und $x$ eingesetzt worden.]}\endgroup
\noindent\rule{\linewidth}{.4pt}
\[
\begin{aligned}
\delta&=+1\\
\delta'&=+3{,}2503\ [-0{,}0055x]\\
\delta''&=-8{,}0450\ [+0{,}010453x]\\
\delta'''&=+13{,}27952[-0{,}015404x]
\end{aligned}
\]
\noindent\rule{\linewidth}{.4pt}\par\medskip
Für ein kurzsichtiges Auge, dessen Weite $k=147\ \mathrm{mm}$\textsuperscript{[**)}, ist
\begingroup\renewcommand{\thefootnote}{}\footnotetext{[**) In der Handschrift steht der Buchstabe $t$ statt $k$.]}\endgroup
\[
\beta'''\mathbin{.}\delta'''\mathbin{.}147=-1,
\]
\noindent also nach Einsetzen der obigen Werte
\[
x^2-864{,}18x=-24641,
\]
\noindent was für $x$ den Wert liefert:]
\[
x=29{,}6\ \mathrm{mm};
\]
\noindent [der zweite aus der Gleichung sich ergebende Wert von $x$ ist aus physikalischen Gründen zu verwerfen.]
\begin{center}Farbiger Rand\end{center}
\[
\begin{array}{r@{\;}r@{\;}r@{\;}r}
[\alpha'\gamma'f'\ ={}]&+4{,}2503-0{,}103596&=&-0{,}440\\
{}[\alpha''\gamma''f''\ ={}]&-4{,}7947-0{,}059363&=&+0{,}285\\
{}[\alpha'''\gamma'''f'''\ ={}]&+5{,}0859+0{,}075266&=&+0{,}382\\\cline{4-4}
&&&+0{,}227
\end{array}
\]

\clearpage\SourceLeaf{127}\HistoricalHead{122}{L}{PHYSIK.\quad NACHLASS.}
[Gesichtsfeld $g=2\mathbin{.}7500''$]
\[
\begin{aligned}
0{,}227\mathbin{.}7500''&=1702''{,}5\text{[*)}\\
\partial n&=0{,}0108\\
n-1&=0{,}515\\
\frac{1702''{,}5\mathbin{.}0{,}0108}{0{,}515}&=]\ 35''{,}7.
\end{aligned}
\]
\begingroup\renewcommand{\thefootnote}{}\footnotetext{[*) Hier fehlt die Bestimmung der Vergrößerung, die notwendig wäre, um die Berechnung des farbigen Randes zu Ende zu führen; theoretisch soll sie gleich $\delta'''$, hier also gleich 13,28 sein. — In der Handschrift steht $1707''{,}5$ statt $1702''{,}5$, entsprechend war auch das Ergebnis zu ändern.]}\endgroup
\begin{center}\rule{.2\linewidth}{.4pt}\end{center}
\clearpage\endgroup
